\documentclass[11pt]{article}

\usepackage[utf8]{inputenc}
\usepackage[T1]{fontenc}
\usepackage{geometry}
\geometry{margin=1in}
\usepackage{amsmath,amssymb}
\usepackage{booktabs}
\usepackage{graphicx}
\usepackage{hyperref}

\title{Synthetic Validation Update for the Universal Embedding Theorem}
\author{Working empirical note}
\date{\today}

\begin{document}
\maketitle

\paragraph{Scope.}
This note summarizes the current synthetic experiment stack in the
project folder for the Universal Embedding Theorem. It is intentionally narrower than the
full theorem draft: the current focus is on synthetic validation before wiring
in the Polymarket and language-model embedding tracks.

\paragraph{Implemented protocols.}
The current runner covers:
\begin{itemize}
  \item PCA recovery of a known causal subspace;
  \item sparse recovery around the $k \log(d / k)$ sample-complexity regime;
  \item a minimum-norm double-descent proxy for benign overfitting;
  \item appended-noise augmentation of the learned embedding;
  \item superposition with random feature dictionaries and sparse activations.
\end{itemize}

\IfFileExists{generated_results.tex}{
  \input{generated_results.tex}
}{
  \paragraph{Missing generated results.}
  Run \texttt{python3 experiments/render\_empirical\_note.py} from the project
  root to materialize the current tables and figure references.
}

\end{document}
